\chapter{Metodologia}

Capitolul de față analizează impactul valorilor lipsă asupra predicției seriilor de timp cu tehnici de învățare automată. Metodologia urmărește o evaluare controlată a modului în care datele incomplete influențează performanța modelelor recurente și compară diferite abordări de tratare a acestora într-un cadru experimental unitar. Sunt acoperite fazele de pregătire și prelucrare a datelor, definirea arhitecturilor analizate, protocolul de antrenare și metricile de evaluare.

\section{Etapele de lucru}

Metodologia experimentală urmează un flux de lucru clar. Seturile de date sunt selectate și încărcate, seriile numerice relevante sunt identificate, observațiile sunt ordonate cronologic și formatate pentru modelele recurente, după care se introduce un mecanism controlat de simulare a valorilor lipsă, astfel încât efectele acestora să poată fi studiate în condiții reproductibile.

După introducerea valorilor lipsă, acestea sunt tratate prin diferite abordări, datele sunt normalizate și împărțite în subseturi de antrenare, validare și testare. Modelele recurente RNN, LSTM și GRU sunt antrenate cu configurații de hiperparametri comparabile, iar predicțiile sunt evaluate prin metrici standard. Rezultatele sunt comparate pentru a măsura efectul valorilor lipsă, al metodei de imputare și al arhitecturii alese.

În paralel cu acest flux principal, sunt rulate și două linii experimentale suplimentare care extind analiza comparativă: învățarea federată și experimentul de ansamblu. Prima urmărește antrenarea unor modele locale pe partiții de date separate și agregarea lor într-un model global, iar a doua urmărește combinarea sau compararea unor predicții obținute din strategii diferite de ansamblare.

\section{Seturile de date utilizate}

Metodologia este validată experimental pe două seturi de date. Primul este DSMTS, un set multivariabil derivat din sisteme dinamice, potrivit pentru studierea dependențelor temporale și pentru experimentarea controlată a strategiilor de injectare și tratare a valorilor lipsă. Al doilea este \textit{Bitcoin Historical Data}, ales pentru volatilitatea ridicată și dinamica temporală complexă pe care le aduce în analiză. Împreună, cele două seturi permit verificarea robusteții concluziilor pe serii de timp cu proprietăți diferite.

\section{Preprocesarea datelor}

Datele sunt preprocesate înainte de antrenarea modelelor pentru a garanta coerența experimentelor și comparabilitatea rezultatelor. Se selectează coloanele relevante, se elimină valorile invalide, se ordonează cronologic observațiile și se extrag secvențele de intrare pentru modelele recurente. Seria de timp este scalată în funcție de experiment: prin standardizare Min-Max, standardizare Z-score sau fără transformare, păstrând valorile originale.

Modelele recurente sunt sensibile la amplitudinea și distribuția valorilor de intrare, motiv pentru care scalarea nu poate fi omisă. Totodată, uniformizarea preprocesării face ca diferențele de performanță observate să poată fi atribuite strategiilor de imputare și arhitecturilor testate, nu variațiilor în reprezentarea datelor.

\section{Simularea și tratarea valorilor lipsă}

Valorile lipsă sunt introduse artificial printr-un mecanism de tip \textbf{MCAR} (\textit{Missing Completely At Random}), pentru a studia efectele datelor incomplete într-un cadru controlat~\cite{rubin1976inference, little2002missing}. Concret, anumite observații din serie sunt eliminate aleatoriu, indiferent de valoarea lor sau de alte variabile disponibile. Această abordare asigură reproductibilitatea și comparabilitatea metodelor de tratare a datelor lipsă între experimente.

După inserarea valorilor lipsă, acestea sunt completate prin patru metode: transmiterea valorii ultime observate, completarea cu zero, completarea cu media coloanei și imputarea predictivă bazată pe \textit{Random Forest}~\cite{breiman2001rf, stekhoven2012missforest}. Primele trei sunt simple și ușor de implementat; ultima estimează valorile lipsă pe baza structurii datelor existente. Compararea acestor strategii arată în ce măsură etapa de preprocesare influențează performanța modelelor de predicție.

\section{Scenariul de învățare federată}

Pentru componenta de învățare federată, seria de timp este împărțită cronologic între mai mulți utilizatori, astfel încât fiecare utilizator primește un subset propriu de antrenare și testare. Fiecare model local este antrenat separat, iar la nivel de server se realizează agregarea rezultatelor într-un model global comun. În implementarea folosită, această procedură este rulată pe mai multe runde federate, iar după fiecare rundă sunt colectate metrici pentru fiecare utilizator și pentru modelul agregat.

Protocolul permite urmărirea atât a evoluției modelelor locale, cât și a efectului pe care agregarea îl are asupra performanței globale~\cite{mcmahan2017fedavg, kairouz2021fl}. Separarea pe utilizatori oferă și posibilitatea de a analiza dacă un anumit mecanism de imputare sau de normalizare se comportă diferit într-un context distribuit. Prin urmare, în lucrarea de față învățarea federată are un dublu rol: testarea unui scenariu de antrenare distribuit și evaluarea robusteții modelelor pe date partiționate.

\section{Experimentul de ansamblu}

În experimentul de ansamblu, sunt testate două topologii: o variantă independentă, în care mai multe modele contribuie separat la agregarea finală, și o variantă secvențială, în care un model corectează sau completează reziduurile altuia. Alegerea acestor două topologii permite o comparație între o strategie de combinare paralelă și una de corecție în lanț.

În ambele variante, evaluarea se face pe aceleași seturi de date și cu aceleași criterii de performanță, astfel încât diferențele observate să poată fi atribuite modului de combinare, nu altor factori experimentali. Ansamblurile sunt sensibile la diversitatea membrilor și la ordinea în care sunt combinate predicțiile~\cite{dietterich2000ensemble, zhou2012ensemble}, ceea ce face această separare necesară. Experimentul completează analiza principală și oferă date suplimentare despre modul în care estimările multiple pot fi combinate pentru rezultate mai stabile.

\section{Arhitectura utilizată}

Lucrarea examinează trei arhitecturi recurente: RNN, LSTM și GRU. Toate aceste modele sunt folosite pentru a face o predicție un pas înainte, pornind de la secvențe fixe de lungime extrase dintr-o serie de timp. Fiecare model primește o fereastră temporală cu observațiile anterioare ca intrare și emite o estimare a valorii următoare. Această formulare face posibilă compararea directă a capacității fiecărei arhitecturi de a modela dependențele temporale în prezența datelor incomplete.

Configurația experimentală permite controlul hiperparametrilor: lungimea secvenței de intrare, dimensiunea stratului ascuns, numărul de straturi recurente, funcția de activare, rata de scădere, rata de învățare, tipul optimizatorului și funcția de pierdere. Pentru anumite scenarii se realizează și analize de tip \textit{hyperparameter sweep} și studii privind impactul normalizării, extinzând analiza dincolo de simpla comparație între arhitecturi și spre sensibilitatea fiecăreia la setările de antrenare.
\section{Protocolul experimental}

Seriile de timp sunt împărțite cronologic în subseturi de antrenare, validare și testare, menținând astfel natura temporală a datelor și evitând utilizarea informațiilor viitoare în procesul de antrenare. Setul de antrenare ajustează parametrii modelului, cel de validare urmărește comportamentul în timpul învățării și ghidează alegerea configurației, iar cel de testare este rezervat evaluării finale.

Metoda necesită antrenamente repetate pentru diferite combinații de arhitectură, strategie de imputare și metodă de normalizare. Această structură permite compararea rezultatelor și izolarea efectului fiecărui factor. Repetarea experimentelor în condiții controlate crește rigorozitatea metodologică și susține formularea unor concluzii mai credibile.

\section{Metricii de evaluare}

Performanța modelelor este evaluată prin metrici uzuale în prognoza seriilor de timp: MSE (\textit{Mean Squared Error}), MAE (\textit{Mean Absolute Error}), RMSE (\textit{Root Mean Squared Error}) și MAPE (\textit{Mean Absolute Percentage Error})~\cite{hyndman2021fpp3, makridakis2018m4}. Utilizarea simultană a acestor metrici acoperă mai multe dimensiuni ale erorii de predicție și oferă o imagine mai completă a calității modelelor.

MSE penalizează mai accentuat erorile mari, RMSE redă eroarea în aceeași unitate cu datele originale, MAE oferă o măsură directă a abaterii absolute, iar MAPE exprimă eroarea în termeni procentuali, facilitând compararea între seturi de date cu scări diferite. Raportarea tuturor acestor valori permite o interpretare mai echilibrată a rezultatelor experimentale.

\section{Sinteza metodologiei}

Cadrul experimental propus în această lucrare este conceput pentru a analiza relația dintre valorile lipsă, strategiile de imputare și performanța modelelor recurente. Combinarea mai multor seturi de date, a unor tehnici de tratare a datelor incomplete cu complexitate variată și a trei arhitecturi neuronale recurente oferă o bază pentru compararea sistematică a rezultatelor.

Experimentele permit măsurarea impactului valorilor lipsă asupra predicției și identificarea strategiilor de tratare a datelor incomplete mai potrivite pentru seriile de timp studiate. Analiza are atât relevanță teoretică, cât și aplicabilitate practică, fiind ancorată în experimente concrete și cuantificabile.
